\documentclass{amsart}
\input{C:/Users/jzchr/MathDocuments/preamble_general.tex}

\title{318 HW 4}
\author{Jalen Chrysos}

\begin{document}
	
\maketitle

\textbf{NOTE TO ALICIA}: The professor gave us an opportunity to redo this problem from last week. This is my same solution from last week because I think it's right (and the other TA marked it correct). I wasn't sure whether resubmitting this was optional or not; if it was optional, please ignore this problem.\\

\textbf{Problem 3.10}: Let $V,W$ be the infinitesimal generators of flows $H,I$. 
\begin{enumerate}[(a)]
	\item Prove that 
	$$
	\partial_s \partial_t\big|_{s=t=0} H_sI_tH_s^{-1}I_t^{-1} = -[V,W].
	$$
	\item Prove that if $H_sI_t=I_tH_s$ for all $s,t\in \R$, then $[V,W]=0$.
	\item Assuming $[V,W]=0$, show that $H_s^*W=W$ for all $s$. Then show that $H_sI_t=I_tH_s$ for all $s,t$.
\end{enumerate}

\begin{proof}
	(a): Viewing $V,W$ as derivations, we can ask how the LHS acts on a function $f:M\to \R$ and expect that we will get $W(V(f))-V(W(f))$:
	\begin{align*}
		\partial_s\partial_t \big|_{s=t=0} H_sI_tH_s^{-1}I_t^{-1} &= \partial_s\partial_t \big|_{s=t=0} H_s^*I_t^*H_s^{-1*}I_t^{-1*}\\
		&= \partial_s\partial_t \big|_{s=t=0} H_s^*I_t^*H_{s*}I_{t*}\\
		&=  \partial_s\big|_{s=0} H_s^*I_0^*H_{s*} (-W) + H_s^*(W)H_{s*}I_{0*}\\
		&= \partial_s\big|_{s=0} (-W) + H_s^*(W)H_{s*}\\
		&= 0 + H_0^*(W)(-V) + (V)(W)H_{0*}\\
		&= -WV + VW\\
		&= -[V,W].
	\end{align*}
	This uses a couple of facts that I should explain: one is that $\partial_t|_{t=0} I_t^*=W$ and $\partial_t|_{t=0}I_{t*}=-W$, and similarly for $s,H,V$ (this is shown in proposition 3.10 in the notes). Another is that when two operators $A_t,B_t$ both depending on a parameter $t$ are composed, their derivative in $t$ works like the product rule:
	$$
	\partial_t|_{t=0} A_t\circ B_t = B_0'\cdot A_0 + A'_0\cdot B_0
	$$
	which follows from treating the two $t$s as separate variables and applying the chain rule.\\
	
	(b): If $H_sI_t=I_tH_s$ then $H_sI_tH_s^{-1}I_t^{-1}=\id$, so by part (a), 
	$$-[V,W] = \partial_s\partial_t H_sI_tH_s^{-1}I_t^{-1} = \partial_s\partial_t \id = 0.$$
	
	(c): By proposition 3.10 in the notes, $\partial_s|_{s=0} H_s^*W = [V,W]$, so if $[V,W]=0$ then $H_s^*W$ is constant in $s$. And at $s=0$ we know that $H_0^*W=W$, so $H_s^*W=W$ for all $s$. The same reasoning shows that $I_t^*V=V$ for all $t\in \R$.
	
	Using this fact, 
	$$H_s^*I_t^*(f) = WV(f) \;\; \text{and} \;\; I_t^*H_s^*(f)=VW(f)$$ and since $[V,W]=0$, $WV(f)=VW(f)$, so $H_s^*I_t^*$ and $I_t^*H_s^*$ act the same way as operators on $\EE(M)$, thus they are equal as derivations and vector fields.
	
	%		 By exercise 3.3, if $V$ generates flow $H$, $H_{t*}(V)=V$ for all $t\in \R$.
\end{proof}

\newpage 

\textbf{Problem 1.2}: Prove that $\gamma_m^1$ is naturally a sub-bundle of the trivial bundle over $\P^m$ of rank $m+1$.

\begin{proof}
	Let $t:E\to \P^m$ be the trivial bundle of rank $m+1$, so $E=\P^m\times \R^{m+1}$, and $\gamma_m^1:E'\to \P^m$ be the tautological line bundle, with 
	$$
	E' := \{(\ell,x)\in \P^m\times \R^{m+1} : x\in \ell\}.
	$$
	Cover $\P^m$ by the open sets $U_j = \{[x_1:x_2:\cdots:x_{m+1}] \in \P^m : x_j\neq 0\}$ for $1\leq j \leq m+1$. For each $U_j$, we have an associated trivialization $\rho_j:E_{U_j}\to U_j\times\R^{m+1}$ given by
	$$
	\rho_j: (\ell,x)\mapsto (\ell,K_{\ell}(x))
	$$ 
	where $K_{\ell}\in \GL_{m+1}(\R)$ is the rotation bringing $\ell$ to $e_j$ (where $e_j$ is the $j$th basis vector of $\R^{m+1}$) via the shortest arc between the two (which is unique if $\ell \in U_j$). If $x\in \ell$, $K_{\ell}(x)\in e_j\R$, so $\rho_j(E'_{U_j})=U_j \times e_j\R$. Thus, this atlas $(U_j,\rho_j)$ makes $\gamma_m^1$ a sub-bundle of $t$.
\end{proof}
\newpage 

\textbf{Problem 1.3}: Prove that a linear form on $\R^{m+1}$ defines a section of the dual bundle of $\gamma_m^1$. Show also that, conversely, any section of this dual determines a function $f:\R^{m+1}\setminus \{0\} \to \R$ with the property that $f(\lambda x)=\lambda f(x)$ for $x\in \R^{m+1}\setminus \{0\}$ and $\lambda \neq 0$.
\begin{proof}
	Let $f:\R^{m+1}\to \R$ be linear, and $\gamma_m^{1*}:E^*\to \P^m$ be the dual bundle of $\gamma_m^1$, so that 
	$$
	E^* := \{(\ell,x^*) : \ell\in \P^m, x^*\in \Hom(\R^{m+1}\cap \ell, \R)\}.
	$$
	Then we can define a section $s:\P^m\to E^*$ by 
	$$
	s: \ell\mapsto (\ell, f|_{\ell})
	$$
	using the fact that $f|_{\ell}$ is linear because $f$ is.\\
%	Let $f:\R^{m+1}\to \big(\bigwedge^k \R^{m+1}\big)^*$ be a linear $k$-form on $\R^{m+1}$. Let $\hat{x}=x/|x|$ for all $x\in \R^{m+1}\setminus \{0\}$. Due to the assumption that $f$ is linear, we know that $f(x)=|x|\cdot f(\hat{x})$.
%	
%	Now define a section of $\gamma_m^{1*}$ by 
%	$$
%	s:\P^m \to E^*, \;\;\;\; s([x^1:\dots:x^{m+1}]) := (y,\ell)\mapsto \<\hat{x},y\> \cdot f(\hat{x})
%	$$
%	where $\<\bullet,\bullet\>$ is the inner product on $\R^{m+1}$. (?)\\
	
	Conversely, let $s:\P^m\to E^*$ be a section of $\gamma_m^{1*}$, so that $\gamma_m^{1*}(s(\ell))=\ell$ for all $\ell \in \P^m$. Then $s(\ell)=(\ell,s_2(\ell))$ where $s_2(\ell)$ is a linear map $\R^{m+1}\cap \ell \to \R$. We can define $f:\R^{m+1}\setminus \{0\}\to \R$ by
	$$
	f(\vec{x}) := s_2([x_1:\cdots:x_{m+1}])(\vec{x}).
	$$
	Note that 
	\begin{align*}
	f(\lambda \vec{x}) &= s_2([\lambda x_1:\cdots: \lambda x_{m+1}])(\lambda \vec{x}) \\
	&= s_2([x_1:\cdots:x_{m+1}])(\lambda \vec{x})\\
	&= \lambda s_2([x_1:\cdots:x_{m+1}])(\vec{x})
	\end{align*}
	because $s_2(\ell)$ is linear.
\end{proof}

\newpage 

\textbf{Problem 1.7}: Let $M,N$ be manifolds and $Q\subset N$ a submanifold. Show that if $f:M\to N$ is a smooth map transversal to $Q$, then the normal bundle of $f^{-1}Q$ in $M$ is (isomorphic to) the pullback under $f:f^{-1}Q\to Q$ of the normal bundle of $Q$ in $N$.
\begin{proof}
	Let $P:= f^{-1}Q$. We know that $P$ is a submanifold of $M$ because $f\pitchfork Q$, and also that $\dim(M)-\dim(P)=\dim(N)-\dim(Q)$, which implies that $\dim(T_pM/T_pP)=\dim(T_qN/T_qQ)$.\\
	
%	At every point $p\in P$, with $q:=f(p)$, 
%	we have the following diagram. The projection $\pi_M$ corresponds to the bundle $\nu(M,P)$, and the projection $\pi_N$ corresponds to $\nu(N,Q)$:
%	$$
%	\begin{tikzcd}
%		{T_pM/T_pP} && {T_qN/T_qQ} \\
%		\\
%		P && Q
%		\arrow["{\pi_M}"', from=1-1, to=3-1]
%		\arrow["{\pi_N}", from=1-3, to=3-3]
%		\arrow["f", from=3-1, to=3-3]
%	\end{tikzcd}
%	$$
%	The fact that the codimension of $P$ is $\dim(N)-\dim(Q)$ implies that $$
%	\dim(T_pM/T_pP) = \dim(M)-\dim(P) = \dim(N)-\dim(Q) = \dim(T_qN/T_qQ)
%	$$
%	so the fibers at $p$ and $q$ are vector spaces of the same dimension, say $r$.\\ 
%	
%	Now consider the pullback $f^*\nu(Q,N)$. The total space is 
%	$$
%	E := \{(x,p)\in T_qN/T_qQ\times P : x\in \pi_N^{-1}(q)\}
%	$$
%	and the map to $P$ is just a projection in the second coordinate. We'll exhibit a vector-bundle isomorphism, i.e. a smooth map $\ph:TM|_P/TP\to E$ that is linear on fibers.
	The normal bundle of $P$ in $M$ is a map $\pi_1:E\to P$ where
	$$
	E := \{(p,v): p\in P, v\in T_pM/T_pP\}, \;\;\; \pi_1:(p,v)\mapsto p
	$$
	and the pullback of the normal bundle of $Q$ in $N$ is $\pi_2:F\to P$ where
	$$
	F := \{(p,w): p\in P, w\in T_qN/T_qQ\}, \;\;\; \pi_2:(p,w)\mapsto p
	$$
	(where $q=f(p)$). To show that these are isomorphic, we must exhibit an isomorphism $\ph:E\to F$ which is linear on fibers. This can be given by
	$$
	\ph: (p,v) \mapsto (q,(D_pf)(v)).
	$$
	It is valid to define $D_pf$ on the quotient $T_pM/T_pP\to T_qN/T_qQ$, since $(D_pf)(T_pP)=T_qQ$. Fixing $p$, $\ph$ is linear because $D_pf$ is. Moreover, restricted to $T_pM/T_pP$, $D_pf$ is invertible as well, since $\dim(T_pM/T_pP)=\dim(T_qN/T_qQ)$ and $D_pf$ is surjective onto $T_qN/T_qQ$ by transversality. 
	
	Thus, $\ph$ is a vector bundle isomorphism between $\nu(P,M)$ and $f^*\nu(Q,N)$, as desired.
\end{proof}


\end{document}